%==============================================================================
% Appendix C: Bianchi I lite comparison and reduced-atlas scope
%==============================================================================
\section{Bianchi I lite comparison and reduced-atlas scope}
\label{app:C}

\subsection{Purpose}
\label{app:C:purpose}

The purpose of this appendix is to make explicit the scope of the present
paper by clarifying the structural differences between the reduced vacuum
atlas of this paper and Bianchi~I / Nieh--Yan-type systems.

\subsection{Structural comparison}
\label{app:C:comparison}

The reduced orbit equation of this paper is $\dot q^{2}+\chi=0$, which
contains only the isotropic scale $q$ and the topology scalar $\chi$.  It
contains neither matter coupling nor effective-source terms.  In contrast,
the Bianchi~I / NY-type bounce-like
systems~\cite{BrandenbergerPeter2017} typically involve several
independent scale factors, matter sources, and effective potentials, and
the existence of a turning point ($\dot q=0$, $\ddot q>0$) is a necessary
condition for the realisation of a bounce-like orbit.

\begin{table}[H]
\centering
\resizebox{\linewidth}{!}{%
\begin{tabular}{lll}
\toprule
Item & Present reduced orbit atlas & Bianchi~I / NY type \\
\midrule
Orbit equation & $\dot q^{2}+\chi=0$ & involves a richer effective potential \\
Scale factor & isotropic scalar $q$ & generally several independent factors \\
Matter coupling & none & generally present \\
Bounce-like turning point & not found in the present scope & depends on the potential \\
Geometric basis & Thurston geometry $\times$ EC+NY torsion & model-dependent \\
\bottomrule
\end{tabular}%
}
\caption{Structural comparison with Bianchi~I lite.}
\label{tab:appC_bianchi}
\end{table}

This comparison is a sanity benchmark, not a matching of the same model.

\subsection{Scope statement}
\label{app:C:scope}

The absence of \code{BOUNCE_LIKE} and \code{RECOLLAPSE_LIKE} orbit classes
in the vacuum reduced atlas of the present paper is a property of this
isotropic vacuum reduced setting.  This result is not a general no-go
claim against matter-coupled, effective-source, or anisotropic extensions.
For instance, on $\Sthree$ ($\chi>0$) the absence of real initial data on
the vacuum reduced atlas is not a denial that the introduction of a matter
source -- changing the structure of the Hamiltonian constraint -- could
permit bounce-like orbits~\cite{BrandenbergerPeter2017}.  The analysis of
these extensions is a future task.
